\begin{afterword}
\summaryhead

The essay defines evidence as an event linked by cause and effect to whatever you want to know about, so that it happens differently depending on how things are. Light bouncing off untied shoelaces into your eyes is evidence about the shoelaces; a machine that beeps for every lottery ticket is evidence about nothing. In technical terms, evidence must share information with its target.

Eyes and brains that work correctly pass this link along, so a rational belief is itself evidence, and when an honest person reports it, it becomes evidence for the listener too. Rational beliefs should therefore spread among honest people. A belief that its holder says cannot be shared, or that is meant to be held whatever happens, is disconnected from reality, and its holder should give it up. The essay ends by asking readers to explain why their own ways of thinking produce true beliefs.

\responsehead

The title asks what evidence is. The essay answers in its second paragraph and then answers twice more, differently. The first answer is physical: a causal link between event and target. The third is about the reader's state of mind: mutual information ``relative to your current state of uncertainty.'' The essay treats these as one idea stated with rising precision. They are two ideas, and the gap between them, between what the world does and what a person is entitled to believe, is a central problem of epistemology. The essay does not address it.

The rest of the essay is not about evidence. It is an argument against private belief, and the argument fails at its central step. Beliefs are said to be ``contagious, among honest folk who believe each other to be honest.'' That makes their spread depend on the listener's trust. A belief that does not spread may lack trust, not truth. The essay drops this condition in the next sentence and calls private reasons ``suspicious.'' It then equates a belief that is not entangled with one that is ``not accurate,'' which a stopped clock refutes, and borrows Tarski's equivalence to finish the job. Tarski, in the paper the epigraph comes from, said his definition of truth ``implies nothing'' about when a sentence may be asserted.

The essay's aim is plainer than its argument. It names a group, ``rationalists,'' who hold a principle it calls ``paradoxical-seeming.'' The principle is fallibilism, and the name is Peirce's. It then names an opponent that is never identified: ``some belief systems,'' which adopt their rules ``in a rather obvious trick.'' No believer's position is stated, so none is answered. The reader receives a test to apply to other people: if your reasons do not persuade me, you should stop believing. That test would also tell an innocent defendant, whose memory the jury does not share, to doubt their own innocence.

The essay ends by asking readers to explain why their thinking produces true beliefs. It never gives that explanation for its own.

In short: the title question gets three wordings of two different definitions, and the rest of the essay turns a failure to persuade into a reason to stop believing, for use against believers it never names.
\end{afterword}
